\documentclass[11pt,a4paper]{article}
\usepackage[utf8]{inputenc}
\usepackage[margin=1in]{geometry}
\usepackage{amsmath,amssymb,amsthm}
\usepackage{listings}
\usepackage{xcolor}
\usepackage{hyperref}

\newtheorem{theorem}{Theorem}
\newtheorem{lemma}[theorem]{Lemma}

\lstset{
  language=C++,
  basicstyle=\ttfamily\small,
  keywordstyle=\color{blue}\bfseries,
  commentstyle=\color{green!60!black},
  stringstyle=\color{red!70!black},
  numbers=left,
  numberstyle=\tiny\color{gray},
  breaklines=true,
  frame=single,
  tabsize=4,
  showstringspaces=false
}

\title{IOI 2010 -- Hotter Colder}
\author{Solution}
\date{}

\begin{document}
\maketitle

\section{Problem Summary}
An interactive problem on the integer line $[1, N]$. A hidden value $X$ is fixed. Starting from position $1$, each guess $G$ receives a response:
\begin{itemize}
    \item \texttt{Hotter} ($+1$): $|G - X| < |P - X|$ where $P$ is the previous guess.
    \item \texttt{Colder} ($-1$): $|G - X| > |P - X|$.
    \item \texttt{Same} ($0$): $|G - X| = |P - X|$.
\end{itemize}
Find $X$ using at most $\lceil \log_2 N \rceil + 1$ queries.

\section{Solution}

\subsection{Key Observation}
If the previous guess is $P$ and the new guess is $G$, the perpendicular bisector of $P$ and $G$ is at position $M = (P + G) / 2$. The response partitions the line:
\begin{itemize}
    \item \texttt{Hotter}: $X$ is on the same side of $M$ as $G$.
    \item \texttt{Colder}: $X$ is on the same side of $M$ as $P$.
    \item \texttt{Same}: $X = M$ (only possible when $P + G$ is even).
\end{itemize}

\subsection{Algorithm}
Maintain a feasible interval $[\mathit{lo}, \mathit{hi}]$ for $X$, initially $[1, N]$. At each step:
\begin{enumerate}
    \item Compute the midpoint $\mathit{mid} = \lfloor (\mathit{lo} + \mathit{hi}) / 2 \rfloor$.
    \item Choose $G = 2 \cdot \mathit{mid} + 1 - P$ so that the bisector of $P$ and $G$ falls at $\mathit{mid} + 0.5$ (between $\mathit{mid}$ and $\mathit{mid}+1$). Clamp $G$ to $[1, N]$.
    \item If \texttt{Hotter} and $G > P$: set $\mathit{lo} = \mathit{mid} + 1$.
    \item If \texttt{Hotter} and $G < P$: set $\mathit{hi} = \mathit{mid}$.
    \item \texttt{Colder} is symmetric.
    \item If \texttt{Same}: $X = (P + G) / 2$ (return immediately).
    \item Update $P \leftarrow G$.
\end{enumerate}

\subsection{Correctness and Query Bound}

\begin{theorem}
The algorithm finds $X$ in at most $\lceil \log_2 N \rceil + 1$ queries.
\end{theorem}
\begin{proof}
After the first query, we have $P \neq 1$ in general, and the feasible interval has been halved. Each subsequent query halves the interval, requiring $\lceil \log_2 N \rceil$ halvings. Adding the initial positioning query gives at most $\lceil \log_2 N \rceil + 1$ total queries.
\end{proof}

\subsection{Complexity}
\begin{itemize}
    \item \textbf{Queries:} $O(\log N)$.
    \item \textbf{Time and space:} $O(1)$.
\end{itemize}

\section{Implementation}

\begin{lstlisting}
#include <bits/stdc++.h>
using namespace std;

// Grader: returns +1 (Hotter), -1 (Colder), 0 (Same)
int Guess(int G);

int HC(int N) {
    int lo = 1, hi = N;
    int prev = 1;

    while (lo < hi) {
        int mid = (lo + hi) / 2;
        int G;
        if (prev <= mid) {
            G = 2 * mid + 1 - prev;
        } else {
            G = 2 * mid - prev;
        }
        G = max(1, min(N, G));

        int result = Guess(G);

        if (result == 0) {
            return (prev + G) / 2;
        }

        int boundary = (prev + G) / 2;
        if (result == 1) {
            // X closer to G
            if (G > prev) lo = boundary + 1;
            else          hi = boundary;
        } else {
            // X closer to prev
            if (G > prev) hi = boundary;
            else          lo = boundary + 1;
        }

        prev = G;
    }

    return lo;
}
\end{lstlisting}

\end{document}
